\documentclass[11pt,a4paper]{ctexart}
\usepackage[margin=2.2cm]{geometry}
\usepackage{booktabs,tabularx,longtable,array,enumitem,xcolor,hyperref,xurl}
\hypersetup{colorlinks=true,linkcolor=blue!55!black,urlcolor=blue!55!black}
\setlength{\emergencystretch}{2em}
\setlist[itemize]{leftmargin=2em,itemsep=0.25em,topsep=0.3em}
\setlist[enumerate]{leftmargin=2.2em,itemsep=0.25em,topsep=0.3em}
\definecolor{light}{RGB}{246,244,239}
\newcommand{\contract}[1]{\par\smallskip\noindent\colorbox{light}{\parbox{0.95\linewidth}{\textbf{本节任务：}#1}}\par\smallskip}
\newcommand{\evidence}[1]{\par\noindent\textbf{所需证据：}#1\par}
\newcommand{\figplan}[1]{\par\noindent\textbf{图表计划：}#1\par}
\newcommand{\boundary}[1]{\par\noindent\textbf{表述边界：}#1\par}
\title{llm-train CC12 Grouped-GEMM Selector 兼容性事件报告提纲\\\large 从 production authority 失败到可复核修复}
\author{事件研究提纲，拟作为独立技术报告或系统论文案例研究}
\date{\today}

\begin{document}
\maketitle

\begin{abstract}
这份报告聚焦 TrainVerify 建立真实双 GPU authority 时确认的一项 llm-train runtime/hardware selector 兼容性问题。固定版本的 llm-train 在 compute capability 12.0 的 RTX PRO 6000 Blackwell 上，把 grouped-GEMM/Triton 选择送入不适用的 B200-only candidate set，导致 production 双 rank compile 路径无法完成。调查将失败收窄到硬件分派条件，并使用一个字节与来源都绑定的最小补丁，让 CC12 回退到 generic A100/H100 Triton autotune candidates。报告同时公开两条撤回记录：cache K/V ownership 的错误归因，以及 balance-loss 语义 bug 的越界推断。它们不是已确认的 llm-train bug。本提纲规定事件时间线、复现、根因、修复、验证和证据边界，避免把兼容性故障扩大成训练语义错误。
\end{abstract}

\section*{报告定位}

\begin{itemize}
  \item \textbf{确认的事件：}llm-train grouped-GEMM/Triton selector 在 CC12 production hardware 上选择了不兼容的候选集合，阻断真实双 rank compile。
  \item \textbf{没有确认的结论：}没有证据表明 grouped-GEMM 的数学算子值错误，也没有证据据此推断模型训练语义错误。
  \item \textbf{分开的另一项 finding：}nnScaler RVD/annotation/adapter 无法表达 post-shuffle zigzag ownership，只完成 CPU/Gloo 双进程源码复现，不能写成 CUDA/NCCL 复现。
  \item \textbf{已撤回：}“cache K/V 是 zigzag-owned”和“balance-loss mask 与 routing map 错位”都不得作为 llm-train bug 结论。
\end{itemize}

\section{执行摘要}
\contract{让维护者不用读调查细节也能知道发生了什么、受什么影响、如何修、验证到了哪一层。}

建议用一页事件卡片：
\begin{longtable}{p{0.20\linewidth}p{0.65\linewidth}}
\toprule
字段 & 内容 \\
\midrule
受影响路径 & llm-train production compile 中的 grouped-GEMM/Triton launch configuration selection \\
触发条件 & compute capability 12.x；真实双 GPU、双 rank authority path \\
用户可见症状 & production authority compile 失败，无法生成可验证 SM/PM 工件 \\
根因位置 & 固定 llm-train revision 的 \texttt{llm/kernel/gemm.py} candidate selector \\
最小修复 & 将两处 \texttt{CC >= 10} 专用分支条件收窄为 \texttt{CC == 10} \\
修复效果 & CC12 使用 generic A100/H100 Triton autotune candidate set \\
不变项 & operator 数学定义、model code、graph annotation、reviewed base revision \\
验证 & 双 rank authority 成功；补丁身份和 patched source hash 进入 receipts/provenance \\
残余风险 & 尚需独立性能回归、更多 CC12 设备和上游维护者确认 \\
\bottomrule
\end{longtable}

\boundary{标题、摘要和结论统一使用“runtime/hardware selector 兼容性问题”。不得简写成“TrainVerify 证明 llm-train 训练结果错误”。}

\section{系统背景}
\contract{解释为什么一个 kernel selector 会出现在形式化验证 authority 链里，以及普通 smoke test 为什么看不到它。}

\subsection{llm-train、nnScaler 与 TrainVerify 的边界}
llm-train 提供 YOCO 模型和 grouped-GEMM kernel 路径；nnScaler trace、规划并生成 SM/PM；TrainVerify 要求从真实 production path 固化 authority 后再翻译为 Lean。用一张三层图说明职责。

\subsection{Grouped GEMM 在 YOCO MoE 中的位置}
用一个 MoE layer 的数据流图解释 router、token dispatch、expert-local up/gate/down projection、grouped GEMM、activation 和 AllToAll。只介绍理解 selector 所需部分，不在事件报告中重讲整个模型。

\subsection{Kernel selector 与数学算子的区别}
说明 operator 定义回答“计算什么”，selector 回答“在当前硬件上使用哪组 launch/autotune candidates”。selector 失败可以阻断编译或运行，但本事件没有显示数学公式发生变化。

\subsection{为什么 CPU smoke 和 planner-only test 不够}
CPU smoke 能检查 import、trace、code generation、pickle extraction 和 shuffle wiring；它不进入真实 CC12 Triton selector、GPU profile 和双 rank NCCL 路径。仅设置 \texttt{plan\_ngpus=2} 也不等于真实启动两张 GPU。

\figplan{图 1：训练源码、nnScaler compile、grouped-GEMM selector、GPU kernel 和 TrainVerify authority 的调用链；用红色标出失败位置。}

\section{环境与 revision 范围}
\contract{把每项结论绑定到实际检查过的源码、硬件和运行环境，防止跨 revision 外推。}

CC12 事件固定在 llm-train \path{9a1be1d5fd1c063d80be82797692cdc7d23cfbef}、nnScaler \path{d3d468ed23edb2f28aa8566b2dfb6ed49c5955cf}，硬件为两张 NVIDIA RTX PRO 6000 Blackwell Workstation Edition，compute capability 12.0，CUDA runtime 12.8，NCCL 2.27.3，driver 595.71.05。模型配置为 YOCO-MoE-A0.4B、24 layers、sequence length 4096。

历史 nnScaler zigzag/RVD finding 使用另一组 revision：nnScaler \path{1102e629ee68ab6f8f4a7c2e721ea894e5962131} 与 llm-train \path{30b80f546d46aacbf8316c983550c50a56bcd1ac}。该 finding 不能自动外推到 CC12 authority 所用版本，也不能与 CC12 root cause 合并。

\figplan{表 1：每条 finding 的 llm-train、nnScaler、TrainVerify revision，硬件路径，复现层级和当前证据状态。}

\section{事件检测}
\contract{按证据发生顺序记录现象，不用事后知识重写成“一眼看出根因”。}

\subsection{预期行为}
固定 llm-train、nnScaler、TrainVerify revision，真实启动两个 rank，在两张 RTX PRO 6000 Blackwell 上完成 production SM/PM compile，并生成带 hardware、profile、rank receipts 和 pickle hashes 的 authority。

\subsection{实际失败}
记录第一条 production \texttt{torchrun} 失败的命令、时间、机器、GPU inventory、软件版本、rank、exit code、完整 stderr 和第一个稳定 stack frame。这里必须引用原始日志；若日志未保存，则明确标记为不可恢复证据，不能补写一个“合理的”报错。

\subsection{为何把它当作 authority BLOCK}
CPU smoke 通过不允许降级替代 production authority。失败发生在真实双 rank compile 路径，因此正确动作是停止 authority 发布并定位硬件 dispatch，而不是手写 PM graph、跳过 profile 或修改 Lean statement。

\evidence{需要找到或明确缺失：首次失败日志、GPU inventory、CUDA、Triton、PyTorch、driver、NCCL 版本、rank receipts，以及失败前后的环境清单。}

\section{诊断时间线}
\contract{把每一步假设、实验、结果和下一步决策连起来。}

版本库能够确认的里程碑为：2026-07-29 修正 nnScaler RVD 归因；2026-07-30 完成 CPU/Gloo 双进程真实源码复现；2026-08-06 落地 CC12 selector 修复与 provenance attestation；2026-08-08 撤回 cache K/V zigzag 归因；2026-08-11 重写公开审计说明。2026-08-06 是修复落地日期，不得冒充首次故障发生日期；首次失败的原始时间仍需从运行日志恢复。

建议同时采用下表记录 CC12 诊断内部顺序，成稿时为 T0 至 T4 补真实时间戳：
\begin{longtable}{p{0.09\linewidth}p{0.16\linewidth}p{0.18\linewidth}p{0.18\linewidth}p{0.16\linewidth}}
\toprule
时间 & 观察 & 假设 & 实验/证据 & 结论 \\
\midrule
T0 & 双 rank production compile 失败 & 图或模型逻辑错误 & 与 CPU smoke、planner-only 对照 & 失败只在真实 GPU path 出现 \\
T1 & authority metadata 显示 CC12 & 硬件 dispatch 不支持 & 固定环境并检查 selector branch & 优先调查 CC 分派 \\
T2 & 失败收敛至 \texttt{gemm.py} & candidate set 与硬件不匹配 & 列出 CC 条件和候选来源 & CC12 进入 B200-only set \\
T3 & 最小条件补丁 & 改 selector 即可恢复 & 只改两处条件，其他源码冻结 & 候选集合改变，算子接口不变 \\
T4 & 双 rank 重跑成功 & 修复充分 & 生成 authority 与 receipts & production path 恢复 \\
\bottomrule
\end{longtable}

每一步都要记录被排除的替代解释，例如缺 profile、nnScaler DP solver、CUDA/NCCL 初始化、模型图错误和 TrainVerify emitter 错误。

\section{根因分析}
\contract{从源码条件、硬件类别和候选集合三方面说明为什么失败，而不是只展示补丁。}

\subsection{触发条件}
给出固定 llm-train revision、文件和两处 selector 条件。解释原条件如何把所有 \texttt{CC >= 10} 设备送入专用分支，而该分支实际对应 B200-only candidate set。

\subsection{为何 CC12 不满足该分支假设}
需要由实际候选定义、Triton/backend约束和原始失败日志共同支持。不能只因为补丁后成功，就反推每个候选为何失败。无法从现有证据确定的低层原因明确留空。

\subsection{五个为什么}
建议按以下链条展开：
\begin{enumerate}
  \item 为什么 production compile 失败？因为 grouped-GEMM selector 没有找到可用的 launch candidate，或在候选编译阶段失败。
  \item 为什么选择了这组 candidates？因为 CC12 满足 \texttt{CC >= 10}。
  \item 为什么这条条件不正确？因为分支能力边界对应 CC10/B200，而不是所有未来更高 compute capability。
  \item 为什么测试没有提前发现？现有测试没有覆盖 CC12 production 双 rank 路径。
  \item 为什么形式化工作触发了它？authority policy 禁止用 CPU 或模拟图替代真实编译，迫使系统进入此前未覆盖的硬件路径。
\end{enumerate}

\boundary{若原始日志只能证明“候选编译失败”，正文就写到这一层；不要补成特定 PTX、寄存器或数值故障。}
\figplan{图 2：原 selector decision tree 与修复后的 decision tree；表 1：CC、设备类别、候选集合和已测试状态。}

\section{最小修复}
\contract{说明补丁为什么最小、改变了什么、没有改变什么，以及如何防止补丁漂移。}

\subsection{补丁内容}
\texttt{scripts/yoco\_regen/patch\_llm\_cc12\_gemm.py} 将 \texttt{llm/kernel/gemm.py} 中两处 \texttt{CC >= 10} 收窄为 \texttt{CC == 10}。补丁标识为：
\begin{quote}
\texttt{cc12\_generic\_triton\_fallback\_v1}
\end{quote}

\subsection{为何选择 fallback 而不是改模型或图}
补丁只改变 launch configuration selection。它不改变 grouped-GEMM 的 tensor contract、模型调用、nnScaler annotation 或 SM/PM topology，因此不会通过“修改被验证程序的数学含义”来绕开失败。

\subsection{补丁完整性}
说明精确替换数、原始字节预期、幂等性、半补丁状态拒绝、patched-file SHA-256 和 rank receipt 绑定。正式运行从 fixed revision 的 private materialization 产生 patched source，不修改开发 worktree。

\section{验证与回归}
\contract{分别验证补丁脚本、production执行、图产物和性能，避免“跑通一次”被写成完整修复证明。}

\subsection{补丁脚本测试}
仓库已有 \path{test_cc12_gemm_patch_is_exact_and_idempotent} 和 \path{test_cc12_gemm_patch_rejects_partial_source}，验证精确两处替换、幂等性和 partial source 拒绝。成稿前还应补充错误 revision、零命中、marker 被篡改和文件存在额外修改的独立测试。该层只能证明文本补丁机制，不能替代 CC12 GPU runtime 复现。

\subsection{真实双 GPU 验证}
在两张 RTX PRO 6000 Blackwell、双 rank NCCL 上重跑 production authority。报告实际命令、环境、exit code、rank receipts 和 authority completion。

\subsection{产物验证}
确认 SM/PM graph、节点数、raw digest、hardware profile、patched source hash 和补丁标识进入 provenance。若修复前没有可比较的完整 authority，不得声称“补丁前后 graph byte-identical”；只能说明成功产物与预期 canonical graph 一致。

\subsection{性能验证}
需要补测 generic fallback 对编译时间、autotune时间和 grouped-GEMM runtime 的影响，并与可用的 CC12 baseline 对照。当前没有数据时写成 open action，不得说“无性能损失”。

\subsection{跨硬件回归}
至少覆盖 CC10 原分支、CC12 fallback 和较低 CC 的既有分支，确认条件收窄没有把 CC10 从专用路径移走。能够访问的硬件和仅做静态检查的范围要分开报告。

\figplan{表 2：修复验证矩阵；表 3：production authority结果；图 3：compile/runtime性能对比。}

\section{影响评估}
\contract{限定受影响条件、故障类型和用户影响。}

\begin{itemize}
  \item \textbf{直接影响：}在受影响硬件与固定源码路径上，production graph compile/authority generation 被阻断。
  \item \textbf{间接影响：}无法生成真实 authority 时，任何后续 Lean proof 都不能合法绑定该硬件路径。
  \item \textbf{未建立的影响：}没有证据显示已完成训练产生错误模型参数或 grouped-GEMM 数值悄悄错误。
  \item \textbf{版本范围：}只对检查过的 llm-train revision 和实际 selector source 下结论；上游其它版本需单独审计。
\end{itemize}

需要在成稿前补充受影响配置矩阵和上游是否已有等价修复。

\section{两个撤回，以及为什么必须写进报告}
\contract{展示调查如何纠正错误假设，防止未来把候选分析当成已确认bug。}

\subsection{撤回一：cache K/V ownership}
早期看到 decoder stream 在 shuffle 后保持 zigzag，一度猜测 cache K/V 也来自 zigzag-owned hidden state。重新沿 \texttt{model.py} producer 顺序审计后发现，K/V projection 在 \path{wrap_maybe_shuffle(h)} 之前完成；K/V 是普通连续 shard，Q 和 decoder stream 才是 shuffle 后 zigzag。建立在错误假设上的候选 commits 不进入 authority，也不构成 llm-train bug 证据。

\subsection{撤回二：balance-loss 语义 bug}
早期调查一度认为 routing map 与 position-indexed loss mask 错位，但目标 TID 是 leaf output，balance-loss consumer 位于 trace scope 之外。用 scope 外代码推断 scope 内 goal 为假属于方法错误。除非固定 revision、源码位置和可复现实验重新建立证据，不得恢复这项结论。

\subsection{与 nnScaler zigzag/RVD finding 的边界}
另一个已确认 finding 是 nnScaler RVD/annotation/adapter integration 无法表达 zigzag ownership，导致普通 rank-order gather 不能还原原序列。它有双进程 CPU/Gloo 真实源码复现，但没有 CUDA/NCCL专项复现。该问题不应并入 CC12 selector root cause，也不应归因到 llm-train API 使用错误。

\figplan{图 4：三条调查线的证据状态图，分别标记 confirmed、withdrawn 和 separately confirmed。}

\section{形式化验证在这次事件中起了什么作用}
\contract{既不夸大 Lean 直接发现 kernel bug，也不忽略严格 authority policy 的作用。}

TrainVerify 的 Lean theorem 没有直接证明 selector 条件错误。真正产生价值的是验证链对输入真实性的要求：
\begin{itemize}
  \item production authority 必须来自真实双 GPU、双 rank 路径，CPU smoke 不能替代；
  \item hardware、profile、patched source 和 graph bytes 必须进入 provenance；
  \item 失败不能通过手写图、弱化命题或修改 metadata 绕过；
  \item 修复后必须重新生成 authority，再重新 emission 和证明。
\end{itemize}

由此可形成更一般的系统论点：proof-oriented validation 会把“测试环境没有覆盖的编译路径”变成必须执行和记录的前置条件。但这一论点需要更多案例与对照实验，不能仅靠单事件扩大成普遍结论。

\section{预防措施与后续行动}
\contract{把教训转成owner明确、可验收的工程动作。}

\begin{longtable}{p{0.23\linewidth}p{0.17\linewidth}p{0.18\linewidth}p{0.17\linewidth}}
\toprule
行动 & 责任层 & 验收标准 & 当前状态 \\
\midrule
用能力表替代开放式 \texttt{CC >= N} 分派 & llm-train & 未知新 CC fail closed 或使用明确generic path & 待上游设计 \\
增加 CC12 selector 单测 & llm-train & 两处 grouped-GEMM path 都选generic candidates & 待补 \\
增加真实 CC12 双 rank CI & 基础设施 & production compile和最小runtime测试通过 & 受硬件资源约束 \\
记录 candidate rejection 原因 & runtime & failure输出候选、backend与首个原因 & 待补 \\
补充性能回归 & llm-train & 编译与runtime指标可复核 & 待测 \\
保留patch provenance & TrainVerify & source hash、patch ID、rank receipt一致 & 已实现 \\
撤回结论登记 & TrainVerify/docs & 候选revision不能进入production allowlist & 已实现 \\
\bottomrule
\end{longtable}

\section{对 OSDI/SOSP 投稿的作用}
\contract{决定这份材料是独立事件论文、主论文案例，还是 artifact appendix。}

单独的两行 selector 修复不足以构成 OSDI/SOSP 论文。更有研究价值的角度是：严格的 proof/authority pipeline 如何暴露传统 smoke test 未覆盖的跨层错误，并如何用 provenance、防绕过和撤回机制约束诊断。是否能独立成文取决于：
\begin{itemize}
  \item 是否有多个跨模型、编译器和硬件的案例；
  \item 是否能与普通 CI、differential testing 和 compiler self-check 做系统对比；
  \item 是否有结构化诊断方法，而不只是事后复盘；
  \item 是否有量化的检测率、定位成本和误报/撤回分析。
\end{itemize}

在证据不足时，这份报告更适合作为 TrainVerify 主论文的详细 case study 与 artifact appendix，同时单独公开为技术事件报告。

\section{相关工作计划}
围绕 GPU kernel dispatch、architecture capability detection、Triton autotuning、ML compiler testing、translation validation、postmortem methodology 和 proof-producing systems 组织。精确比较必须查阅原始论文与上游 issue/commit，提纲阶段不写未经核验的优先权结论。

\section{结论}
结论保持三句话：发生的是 CC12 grouped-GEMM selector 兼容问题；最小补丁恢复了真实双 rank authority path，并被 provenance 绑定；该事件没有证明 llm-train 的训练数学语义错误。随后指出，production-path evidence 和可撤回的诊断记录是 TrainVerify 方法的重要组成部分。

\appendix
\section{计划附录}
\begin{itemize}
  \item A：完整环境与revision清单。
  \item B：原始失败日志和符号化stack trace。
  \item C：selector源码、最小补丁与精确hash。
  \item D：补丁脚本的 fail-closed 测试，并区分已存在测试与待补测试。
  \item E：双 rank复现命令、rank receipts和authority manifest。
  \item F：性能与跨硬件回归。
  \item G：K/V与balance-loss撤回的证据时间线。
  \item H：nnScaler zigzag/RVD独立finding的复现边界。
\end{itemize}

\section*{成稿前的BLOCK}
\begin{itemize}
  \item 恢复或明确缺失首次production失败的原始日志、时间戳和完整环境。
  \item 用源码和日志证明“B200-only candidate set”与失败之间的确切关系，不从成功补丁反推未记录细节。
  \item 运行补丁前后可比较的性能实验；没有数据时不得声称无性能影响。
  \item 检查llm-train上游当前版本是否已修复，并限定受影响revision范围。
  \item 验证CC10、CC12和较低CC路径，区分真实硬件测试与静态selector测试。
  \item 为所有图、表和结论建立claim-to-evidence ledger。
\end{itemize}

\end{document}
